\section{结论}

本文提出QDPT，通过问题Token池化与视觉Cross-Attention生成语言侧动态Prompt，用于冻结多模态大模型的领域适配。在Qwen3-VL-4B-Instruct上，QDPT使用约7.81M可训练参数，PathVQA Test准确率达到58.88\%，较静态语言Prompt提高2.80个百分点。问题条件Query在三次参数量匹配对照中平均高于Learned Query 0.83个百分点，但优势随随机种子变化；关闭视觉读取或移除静态视觉Prompt也会降低准确率。统一条件下的短程测试测得1.88倍于LoRA-r8的训练吞吐。电气选择题结果补充了设备图像上的应用验证。SLAKE的知识型问答与开放回答、PathVQA的运行稳定性仍是QDPT的主要不足。后续工作将改进问题分支的词序表达，并降低动态生成器对初始化的敏感性。
